\documentclass[11pt]{article}
\usepackage[utf8]{inputenc}
\usepackage{geometry}
\geometry{a4paper, margin=1in}
\usepackage{hyperref}
\usepackage{booktabs}
\usepackage{enumitem}
\usepackage{amsmath}
\usepackage{float}
\usepackage{xcolor}

\title{Lifecycle 3 --- Q\&A Preparation \\ \large Questions Your Guide Will Ask}
\author{Study Reference}
\date{2026-03-18}

\begin{document}

\maketitle

\tableofcontents
\newpage

% ============================================================================
\section{Focal Loss Questions}
% ============================================================================

\subsection{Q: What is Focal Loss? How is it different from Cross-Entropy?}

\textbf{Cross-Entropy:} Treats all samples equally. Easy examples (model is 99\% sure and correct) get the same importance as hard examples (model is 55\% sure).

\textbf{Focal Loss:} Adds a modulating factor that down-weights easy examples:
\[
\text{FL}(p_t) = -(1 - p_t)^\gamma \cdot \log(p_t)
\]

\textbf{Numerical example:}
\begin{itemize}[noitemsep]
    \item Easy example ($p_t = 0.9$): CE = 0.105, FL = 0.00105 $\to$ \textbf{100$\times$ less loss}
    \item Hard example ($p_t = 0.3$): CE = 1.204, FL = 0.590 $\to$ retains significant loss
\end{itemize}

The network ignores easy examples and concentrates learning on hard borderline cases.

\subsection{Q: Why Focal Loss instead of just oversampling?}

Both address class imbalance but differently:

\begin{table}[H]
    \centering
    \renewcommand{\arraystretch}{1.3}
    \begin{tabular}{lp{5cm}p{5cm}}
        \toprule
        & \textbf{Oversampling} & \textbf{Focal Loss} \\
        \midrule
        How & Duplicates minority class images & Changes the loss function \\
        Risk & Overfitting on duplicates & None (no data change) \\
        Focus & Class membership & Sample difficulty \\
        \bottomrule
    \end{tabular}
\end{table}

Focal Loss is smarter because a non-landslide image can also be ``hard'' (e.g., ploughed field that looks like a landslide). Focal Loss focuses on difficulty, not just class.

\textbf{We use both:} WeightedRandomSampler (data-level) + Focal Loss (loss-level).

\subsection{Q: Why $\gamma = 2.0$? Why not higher or lower?}

\begin{table}[H]
    \centering
    \renewcommand{\arraystretch}{1.2}
    \begin{tabular}{lcl}
        \toprule
        $\gamma$ & \textbf{Behavior} & \textbf{When to use} \\
        \midrule
        0 & Standard Cross-Entropy & Balanced datasets \\
        1 & Mild focus on hard examples & Slight imbalance \\
        \textbf{2} & \textbf{Strong focus} & \textbf{Our case: 72/28 imbalance} \\
        5 & Very aggressive & Extreme imbalance ($>$95/5) \\
        \bottomrule
    \end{tabular}
\end{table}

$\gamma = 2.0$ is the default from the original paper (Lin et al., ICCV 2017). For our moderate 72/28 imbalance, it provides strong focus without being too aggressive. A sweep over $\gamma$ values is a listed future improvement.

\subsection{Q: What is the $\alpha$ parameter in Focal Loss?}

$\alpha$ is an optional per-class weighting factor. The full formula is:
\[
\text{FL}(p_t) = -\alpha_t (1 - p_t)^\gamma \cdot \log(p_t)
\]

We did NOT use $\alpha$ because we already handle class imbalance at the data level with WeightedRandomSampler. Using both $\alpha$ and the sampler could over-correct. Testing $\alpha$ is a future improvement.

% ============================================================================
\section{Test-Time Augmentation Questions}
% ============================================================================

\subsection{Q: What is TTA? Why 5 views?}

TTA = Test-Time Augmentation. Instead of one forward pass, we create 5 different views of the same image and average predictions:

\begin{enumerate}[noitemsep]
    \item Original (no change)
    \item Horizontal flip (mirror left-right)
    \item Vertical flip (mirror top-bottom)
    \item Both flips combined
    \item 90\textdegree\ rotation
\end{enumerate}

\textbf{Why 5:} These geometric transforms are guaranteed to preserve landslide features (a landslide flipped is still a landslide). We avoid color/brightness augmentations at test time because they could distort satellite imagery characteristics.

\subsection{Q: Doesn't TTA make inference 5$\times$ slower?}

Yes, but:
\begin{enumerate}[noitemsep]
    \item \textbf{Accuracy $>$ speed} for disaster monitoring. Extra 3 seconds is acceptable.
    \item \textbf{GPU parallelism:} Batch all 5 views $\to$ actual slowdown is $\sim$2$\times$, not 5$\times$
    \item \textbf{Adaptive TTA (future):} Only use multi-view when single-pass confidence is borderline (40--55\%). High-confidence images skip TTA.
\end{enumerate}

\subsection{Q: How does TTA help mathematically?}

Averaging reduces variance. If single-pass has noise $\sigma$, averaging $n$ views has noise $\sigma / \sqrt{n}$. With 5 views, noise is reduced by $\sqrt{5} \approx 2.24\times$.

\textbf{Practical effect:} A borderline case with $P = 0.48$ (miss at threshold 0.467) might become $P_{\text{TTA}} = 0.51$ (correct detection) because the other views provide cleaner signals.

% ============================================================================
\section{Grad-CAM Questions}
% ============================================================================

\subsection{Q: What is Grad-CAM? Can it improve accuracy?}

Grad-CAM is an \textbf{interpretability tool}, NOT an accuracy improvement. It generates a heatmap showing which image regions the model focused on.

\textbf{How it works:}
\begin{enumerate}[noitemsep]
    \item Forward pass: get prediction
    \item Backward pass: compute gradients of predicted class w.r.t.\ last convolutional layer
    \item Gradients become importance weights for each feature channel
    \item Weighted sum of feature maps $\to$ heatmap
    \item ReLU (only positive influence) $\to$ normalize $\to$ overlay on image
\end{enumerate}

\textbf{Formula:}
\[
L_{\text{Grad-CAM}}^{c} = \text{ReLU}\!\left(\sum_k \alpha_k^c \cdot A^k\right), \quad \alpha_k^c = \frac{1}{Z} \sum_i \sum_j \frac{\partial y^c}{\partial A^k_{ij}}
\]

\subsection{Q: Why is Grad-CAM important for this project?}

Because this is a \textbf{safety-critical} system:
\begin{itemize}[noitemsep]
    \item Disaster teams won't issue evacuation orders based on a black-box number
    \item Grad-CAM lets them visually verify: ``Is the model looking at actual debris, or at a cloud shadow?''
    \item If hot regions match geological features $\to$ trust the prediction
    \item If hot regions are on irrelevant areas $\to$ flag as unreliable
    \item Explainability is increasingly required by regulators for AI in safety-critical applications
\end{itemize}

\subsection{Q: Why is Grad-CAM resolution so low (10$\times$10)?}

The heatmap resolution matches the spatial size of the target convolutional layer:
\begin{itemize}[noitemsep]
    \item EfficientNet-B3 final feature map: 10$\times$10 $\to$ coarse heatmap
    \item ViT-B/16: 14$\times$14 patches $\to$ slightly better
    \item The heatmap is upsampled (interpolated) to full image size for overlay
\end{itemize}

\textbf{Improvement:} Grad-CAM++ or Score-CAM produce finer heatmaps. For ViT, attention rollout across all layers gives richer spatial maps. These are listed as future improvements.

% ============================================================================
\section{Results Questions}
% ============================================================================

\subsection{Q: Your r6 metrics are projected, not actual. Why?}

Focal Loss + TTA require full retraining of both models and complete evaluation. The projections are based on:
\begin{itemize}[noitemsep]
    \item Published literature: Focal Loss $\to$ +2--4\% F1 on imbalanced datasets
    \item Published literature: TTA $\to$ +1--3\% accuracy on image classification
\end{itemize}

\textbf{How to answer this confidently:} ``The retraining is the immediate next step. The projected ranges give us a validation target. The methodology --- which is the main thesis contribution --- is fully implemented and ready to run.''

\subsection{Q: Why did multi-class drop to 71.67\%?}

Three reasons:
\begin{enumerate}[noitemsep]
    \item \textbf{Too few per-class samples:} Some classes have very few training images
    \item \textbf{Visual similarity:} Mudflows and debris flows look nearly identical from satellite altitude
    \item \textbf{6-class is harder:} Binary needs 1 boundary, 6-class needs 5 boundaries
\end{enumerate}

\textbf{Solution:} Hierarchical classification --- binary detection first (high accuracy), then type classification only on positives.

\subsection{Q: What is the overall improvement from r0 to r6?}

\begin{table}[H]
    \centering
    \renewcommand{\arraystretch}{1.2}
    \begin{tabular}{lccc}
        \toprule
        \textbf{Metric} & \textbf{r0 (start)} & \textbf{r6 (projected)} & \textbf{Total gain} \\
        \midrule
        Accuracy & 78.54\% & 86--88\% & +7.5--9.5\% \\
        F1-Score & 67.67\% & 78--81\% & +10.3--13.3\% \\
        ROC-AUC & 85.53\% & 92--94\% & +6.5--8.5\% \\
        \bottomrule
    \end{tabular}
\end{table}

% ============================================================================
\section{General Project Questions}
% ============================================================================

\subsection{Q: What is the main contribution of your project?}

\begin{enumerate}[noitemsep]
    \item \textbf{Dual-phase pipeline:} CNN detection + HMM forecasting (not just detection)
    \item \textbf{CNN-Transformer hybrid ensemble:} local + global features
    \item \textbf{Systematic iteration:} 7 experiments, 3 lifecycles, each motivated by identified gaps
    \item \textbf{Explainability:} Grad-CAM for safety-critical trust
\end{enumerate}

\subsection{Q: How is your work different from existing papers?}

\begin{itemize}[noitemsep]
    \item Most papers use single model --- we use architecturally diverse ensemble
    \item Most papers only detect --- we also forecast future risk (HMM)
    \item Most papers don't calibrate confidence --- we use temperature scaling
    \item Most papers don't explain predictions --- we use Grad-CAM
\end{itemize}

\subsection{Q: What would you do differently if starting over?}

\begin{enumerate}[noitemsep]
    \item Use 4-channel input (R, G, B, NIR) for NDVI from the start
    \item Use larger dataset (Landslide4Sense)
    \item Implement Focal Loss from the beginning
    \item Automated hyperparameter search (Optuna)
\end{enumerate}

\subsection{Q: Can this be deployed in real-world monitoring?}

Close but needs:
\begin{enumerate}[noitemsep]
    \item Validate r6 with actual retraining
    \item Test on diverse geographies beyond HR-GLDD
    \item Real-time inference (ONNX/TensorRT)
    \item Pixel-level segmentation (U-Net) for precise boundaries
    \item Integration with satellite feeds (Sentinel-2)
\end{enumerate}

\subsection{Q: What is the novelty of using HMM with CNN?}

Most landslide detection papers are purely spatial --- they classify one image at a time with no temporal context. Our HMM adds:
\begin{itemize}[noitemsep]
    \item \textbf{Type prediction:} What kind of landslide (shallow, deep, debris, rockfall)
    \item \textbf{Occurrence probability:} How likely is this event based on regional history
    \item \textbf{Peak risk month:} When is this region most vulnerable (e.g., July monsoon)
    \item \textbf{Future forecast:} Probability of recurrence in next 1--3 time steps
\end{itemize}

This temporal layer transforms a static ``yes/no'' detector into an actionable risk assessment system.

\end{document}
